\documentclass[11pt]{article}
\usepackage[margin=0.85in]{geometry}
\usepackage{amsmath,amssymb,booktabs,xcolor}
\usepackage[colorlinks=true,linkcolor=blue,urlcolor=blue]{hyperref}
\definecolor{ink}{RGB}{23,65,100}
\newcommand{\ket}[1]{\lvert#1\rangle}
\newcommand{\bra}[1]{\langle#1\rvert}
\newcommand{\Tr}{\operatorname{Tr}}
\newcommand{\D}{\mathcal D}
\newcommand{\lesson}[2]{\section*{\color{ink}#1}\addcontentsline{toc}{section}{#1}\noindent\textit{#2}\par\medskip}
\newcommand{\example}{\par\medskip\noindent\textbf{Worked example.}\ }
\newcommand{\takeaway}[1]{\par\medskip\noindent\textbf{Remember:} #1\par}
\setlength{\parindent}{0pt}
\setlength{\parskip}{6pt}
\setlength{\emergencystretch}{2em}
\begin{document}
\begin{center}
{\LARGE\bfseries\color{ink} Understanding the Nov17 circuit}\\[5pt]
{\large Ten short lessons with equations and worked examples}\\[5pt]
{A companion to \texttt{nov17\_methods.pdf} \quad | \quad November 17, 2026}
\end{center}
\lesson{1. What is actually inside the device?}{Two cavities store six encoded spins. Three physical transmons control them.}
A cavity is an oscillator: its Fock state $\ket n$ contains exactly $n$ photons.
We use eight levels, $n=0,\ldots,7$, as the eight basis states of three encoded
bits. This is a choice of labels; there are not three separate physical
qubits inside a cavity.
\begin{equation}
 \boxed{\ket{b_2b_1b_0}_{\rm encoded}\ \longleftrightarrow\
 \ket n_{\rm cavity},\qquad n=4b_2+2b_1+b_0.}
\end{equation}
\example $\ket{101}$ means $n=4+0+1=5$ photons. Cavity $A$ stores
spins $(0,2,4)$ from the least to the most significant bit; cavity $B$
stores $(1,3,5)$. Thus $\ket{5}_A\ket{2}_B$ represents spin bits
$(s_0,s_1,s_2,s_3,s_4,s_5)=(1,0,0,1,1,0)$.

There are two local control transmons, $q_A$ and $q_B$, plus a shared bus
$q_b$. Each is approximated by two levels, $\ket g$ and $\ket e$.
The coupling graph is
\begin{center}
\begin{tabular}{ccccc}
$q_A$&&&&$q_B$\\
$\vert$&&&&$\vert$\\
Cavity $A$&---&Bus $q_b$&---&Cavity $B$
\end{tabular}
\end{center}
The physical system starts in vacuum with all transmons in $g$.
A local preparation pulse creates the uniform encoded superposition:
\begin{equation}
 \ket0_{\rm cavity}\ \longrightarrow\
 \frac1{\sqrt8}\sum_{n=0}^{7}\ket n
 \ \longleftrightarrow\ \ket+^{\otimes3}.
\end{equation}
\example After this ideal preparation, each of the eight Fock outcomes has
probability $1/8$. Its mean photon number is $(0+1+\cdots+7)/8=3.5$.
Starting from vacuum and later having a plus state are consecutive steps.

\textbf{Notation used throughout.} Lowercase $a,a^\dagger$ act on physical
cavity levels; $n=a^\dagger a$ counts photons. Physical transmons use
$\sigma_z,\sigma_\pm$, with $\sigma_z\ket g=-\ket g$ and
$\sigma_z\ket e=+\ket e$. Encoded spins use capital $X_i,Z_i$, with
$Z\ket0=+\ket0$ and $Z\ket1=-\ket1$. For example, $n\ket5=5\ket5$ and $a\ket5=\sqrt5\ket4$:
the latter lowers the photon number. These are different objects.
\takeaway{The photons store the state. The transmons help change its amplitudes and phases.}

\clearpage
\lesson{2. What state do we want to prepare?}{The spin Hamiltonian defines the optimization goal.}
For six encoded spins, the target model is
\begin{equation}
 \boxed{H_{\rm SK}=\sum_{i<j}J_{ij}Z_iZ_j-B\sum_iX_i.}
\end{equation}
There are $\binom62=15$ pair terms. $J_{ij}$ specifies whether a pair favors
matching or opposite $Z$ signs. The transverse field $B$ favors positive
$X$ polarization. All spin energies here use the project's fixed energy
units, separate from the MHz used for physical pulse controls.

\example Take just two spins with $J_{01}=0.4$ and $B=0.2$.
For $\ket{00}$, the pair product is $+1$, and both $X$ expectations are zero:
its energy is $+0.4$. For $\ket{01}$ the pair product is $-1$, so the
energy is $-0.4$. For $\ket{++}$ the pair expectation is zero and each
$X$ expectation is $+1$, giving energy $-0.2(1+1)=-0.4$.
These candidate energies do not by themselves identify the exact ground state.

The ground state minimizes energy over all normalized states:
\begin{equation}
 H_{\rm SK}\ket{g_{\rm SK}}=E_0\ket{g_{\rm SK}},\qquad
 E(\boldsymbol\theta)=\bra{\psi(\boldsymbol\theta)}H_{\rm SK}
                         \ket{\psi(\boldsymbol\theta)}.
\end{equation}
\example In our saved six-spin instance, with seed $831106000$ and
$B=0.183$, exact diagonalization gives $E_0=-2.7968862$.
The optimized ideal $L=1$ ansatz has $E=-2.6208718$. Its excess energy is
$0.1760144$ in total, or $0.0293357$ per spin. The optimizer tries to reduce
this difference by changing the ansatz parameters.

\textbf{Why introduce another Hamiltonian later?}
$H_{\rm SK}$ says what state we want. The physical cavity--transmon
Hamiltonian says which microwave pulses can prepare it. The experiment
compiles the desired spin operations into physical controls; it does not
acquire a literal direct $Z_iZ_j$ coupling between every pair of encoded bits.

\textbf{Reading the three energy symbols.}
\begin{center}
\begin{tabular}{ll}
\toprule
Symbol & Meaning\\\midrule
$H_{\rm SK}$ & Target spin model\\
$H_{\rm sys}(t)$ & Driven physical device\\
$E(\boldsymbol\theta)$ & Measured or computed objective for the optimizer\\\bottomrule
\end{tabular}
\end{center}
\takeaway{First choose the desired spin evolution; then use the device Hamiltonian to implement it.}

\clearpage
\lesson{3. Which physical Hamiltonian controls the device?}{Fixed dispersive coupling makes frequencies depend on photon number; drives do the active control.}
In rotating frames, use $h=H_{\rm sys}/\hbar$ so that its units are angular
frequency. The total baseline for two cavities and three transmons is
\begin{equation}
\boxed{\begin{aligned}
 h(t)={}&\underbrace{\sum_{j=A,B}\widetilde\Delta_{c,j}n_j
           +\sum_{r=A,B,b}\frac{\Delta_{q,r}}2\sigma_r^z}_{\text{detunings}}\\
 &\underbrace{-\frac12\sum_{r,j}\chi_{rj}n_j\sigma_r^z}_{\text{dispersive coupling}}\\
 &+\underbrace{\sum_j(\epsilon_j a_j^\dagger+\epsilon_j^*a_j)
       +\sum_r(\Omega_r\sigma_r^++\Omega_r^*\sigma_r^-)}_{\text{applied drives}}.
\end{aligned}}
\end{equation}
$\Delta$ means mode frequency minus drive-frame frequency. The four
nonzero shifts are $\chi_{AA}=\chi_A$, $\chi_{BB}=\chi_B$,
$\chi_{bA}$, and $\chi_{bB}$. The complex envelopes $\epsilon$ and $\Omega$
contain the two I/Q drive quadratures: the real and imaginary parts that
set drive strength and phase. Optional Kerr and exchange terms
are zero in the present baseline; the methods PDF specifies their extensions.

\example With local $\chi_A/(2\pi)=1$ MHz, adding one photon lowers the
local transmon's transition frequency by $1$ MHz. Cavity drives mix photon
numbers; transmon drives rotate the physical transmon. The fixed coupling
lets those rotations respond differently to different photon numbers.

\textbf{The factor of two matters.} The $\chi$ quoted here is the full
transmon transition shift per photon. Because the two $\sigma_z$
eigenvalues differ by two, the Hamiltonian coefficient is $-\chi/2$.
To preserve the configured ground-transmon cavity frequencies, write
\begin{equation}
 \widetilde\Delta_{c,j}=\Delta_{c,j}-\frac12\sum_r\chi_{rj}.
\end{equation}
\example For cavity $A$, the local and bus shifts are both $1$ MHz and the
configured $\Delta_{c,A}=0$. Hence $\widetilde\Delta_{c,A}/(2\pi)=-1$ MHz.
With both transmons in $g$, their dispersive contributions give
$+0.5n_A+0.5n_A$ MHz and exactly cancel $-n_A$ MHz. Dropping this cavity
half-shift would change the dynamics.

\textbf{Units and drive convention.} Time is in $\mu$s; multiply MHz by
$2\pi$ to obtain rad/$\mu$s. Our transmon drive has no extra $1/2$:
its resonant angular Rabi frequency is $2|\Omega|$.
For $\Omega/(2\pi)=0.1$ MHz, a resonant $\pi$ rotation takes
$\pi/(2|\Omega|)=2.5\,\mu$s. Dissipation rates $1/T$ have no $2\pi$ factor.
\takeaway{Dispersive terms stay present. A pulse schedule chooses which drives to apply.}

\clearpage
\lesson{4. How does one bus pulse implement joint SNAP?}{It gives each joint photon-number component its own phase.}
During the joint phase gate, drive only the shared bus. In a fixed cavity
component $\ket{n,m}$, the bus sees the conditional detuning
\begin{equation}
 \boxed{\delta_{nm}=\Delta_{q,b}-\chi_{bA}n-\chi_{bB}m.}
\end{equation}
\example Our baseline gives $\delta_{nm}/(2\pi)=8.05-n-1.3m$ MHz.
The components $\ket{1,0}$ and $\ket{0,1}$ therefore correspond to
$7.05$ and $6.75$ MHz. They can acquire different phases under a shaped
bus drive. Across the logical $8\times8$ space, the smallest separation
is only $0.1$ MHz, so selectivity is a real pulse-design constraint.

The desired cavity operation is
\begin{equation}
 \boxed{S_{AB}(\Phi)=\sum_{n,m=0}^{7}e^{-i\Phi_{nm}}
                                \ket{n,m}\bra{n,m}.}
\end{equation}
The bus starts in $g$, moves during the pulse, and should return to $g$
with the desired phase for each component. Our compiler optimizes its I/Q
waveform to approximate this operation simultaneously in all 64 components.

\example Suppose a toy phase table gives only $\ket{1,1}$ a phase $\pi$.
Applied to $(\ket0+\ket1)_A(\ket0+\ket1)_B/2$, it produces
$(\ket{00}+\ket{01}+\ket{10}-\ket{11})/2$. The probabilities remain
$1/4$, but the state is entangled. The coefficient matrix
$\bigl(\begin{smallmatrix}1&1\\1&-1\end{smallmatrix}\bigr)/2$ has rank two,
whereas a product state has rank one.

\textbf{Why a shared drive is necessary.} Independent local phases have
$\Phi_{nm}=\alpha_n+\beta_m$ modulo $2\pi$. For the toy table, the combination
$\Phi_{00}+\Phi_{11}-\Phi_{01}-\Phi_{10}=\pi$ cannot arise from that sum.
The joint bus control provides the required nonseparable phase.
An undriven bus held in $g$ does not implement this arbitrary phase table.

\textbf{How the papers fit.} Heeres et al. introduce number-selective
SNAP phases in one storage cavity; our joint operation extends that idea
using a bus coupled to two storage modes. Wang et al. provide an
experimental shared-transmon/two-cavity precedent. Zhang et al.
(arXiv:1808.07960) analyze driven exchange and beam-splitter interactions,
which move photons between cavities. That is a different operation from
this number-diagonal phase gate.
\takeaway{Joint SNAP changes phases without changing ideal Fock populations. The bus drive and both bus dispersive shifts work together.}

\clearpage
\lesson{5. What does local GRAPE actually do?}{GRAPE chooses microwave waveforms; it is not an additional physical gate or coupling.}
With the bus in $g$ and its drive off, the ideal baseline separates into two
local control problems. For cavity $j$ and its local transmon,
\begin{equation}
\boxed{\begin{aligned}
 h_j(t)={}&(\Delta_{c,j}-\chi_j/2)n_j
       +\frac{\Delta_{q,j}-\chi_j n_j}{2}\sigma_j^z\\
 &+\epsilon_j(t)a_j^\dagger+\epsilon_j^*(t)a_j
       +\Omega_j(t)\sigma_j^++\Omega_j^*(t)\sigma_j^-.
\end{aligned}}
\end{equation}
\example At zero configured detunings and $\chi_j/(2\pi)=1$ MHz, the drift
on $\ket{g,n}$ is zero: $-n/2+n/2=0$. On $\ket{e,n}$ it is $-n$ MHz.
Turning on the cavity drive mixes Fock levels; driving the transmon accesses
this number-dependent evolution. Both controls are available during GRAPE.

Divide the pulse into short constant-amplitude bins. GRAPE adjusts the four
real amplitudes per bin to improve overlap with a target. It uses derivatives
of the objective to decide which amplitudes to change. In the saved
preparation pulse there are 96 bins: about $19.9$ ns per bin over $1.91\,\mu$s.

\textbf{Two different targets use the same controls.}
\begin{equation}
 \begin{aligned}
 \text{Preparation: }&U_{\rm prep}\ket{g,0}
       \simeq\ket g\frac1{\sqrt8}\sum_{n=0}^{7}\ket n,\\
 \text{Coherent Hadamard: }&U_H\ket{g,n}
       \simeq\ket g H^{\otimes3}\ket n\quad\text{for every }n=0,\ldots,7.
 \end{aligned}
\end{equation}
\example Both targets agree on input $n=0$. On input $n=1$, a true
three-bit Hadamard produces alternating signs:
$(\ket0-\ket1+\ket2-\ket3+\ket4-\ket5+\ket6-\ket7)/\sqrt8$.
The preparation-only pulse has no requirement to do that. It can therefore
be easier to optimize and shorter: our selected preparation takes
$1.91\,\mu$s, while the shorter reusable coherent H takes $7.64\,\mu$s.

\textbf{Local shell SNAP.} To apply a local diagonal phase table, turn off
the cavity drive and shape the local-transmon drive. For example, a desired
$\ket5\mapsto e^{-i0.3}\ket5$ can be part of a local SNAP table.
A general H requires population mixing and uses the richer local drive set.

\textbf{Readout versus coherent control.} A measurement pulse may produce
$D H^{\otimes3}$, where $D$ only changes output phases: number probabilities
are unchanged. An interior coherent H must preserve the required relative
phases because later pulses can turn them into probability differences.
\takeaway{Use a state-preparation objective at the start, coherent H inside the circuit, and a measurement objective only at readout.}

\clearpage
\lesson{6. What is the Nov17 ansatz, step by step?}{Each layer applies a joint cost phase and a local mixer expressed through Hadamards.}
Define $W=H^{\otimes6}$, implemented by one three-bit H per cavity in parallel.
The ideal state evolves as
\begin{equation}
 \boxed{\ket{\psi_0}=\ket+^{\otimes6},\qquad
 \ket{\psi_\ell}=W S_\ell^Z W C_Z\ket{\psi_{\ell-1}}.}
\end{equation}
\example For $L=1$, read from right to left: prepare plus, apply $C_Z$,
apply $W$, apply $S_1^Z$, and apply the last $W$. For the common native-X
output, stop before the last $W$. The pulse durations are then
$1.91+10+7.64+2=21.55\,\mu$s. The two local pulses run in parallel,
so their durations are counted once, not twice.

The cost phase uses three matrices:
\begin{equation}
 C_Z=e^{-i\sum_{i<j} A_{ij}Z_iZ_j},\qquad
 A=\gamma(J+aK+bM).
\end{equation}
$J$ is the random coupling matrix. Form $K$ from $J^2$ with its diagonal
removed, then rescale its upper-triangle RMS (root mean square of the
15 distinct pair coefficients) to equal that of $J$.
For example, a path with $J_{01}=0.2$ and $J_{12}=0.3$ contributes
$0.06$ to the unnormalized $K_{02}$. For $M$, fill the nine cross-cavity entries with the mean of the corresponding
$J$ entries, set within-cavity entries to zero, and apply the same RMS
normalization. A zero template stays zero. $K$ is a matrix square, not the
square of the spin Hamiltonian.

\example For the saved $L=1$ instance, $J_{01}=0.0282334$,
$K_{01}=-0.2570185$, and $M_{01}=0.5596066$.
With $(\gamma,a,b)=(0.4635323,-0.2182316,-0.0488731)$,
$A_{01}=0.0264089$ radians. This pair alone gives matching bits
$e^{-i0.0264089}$ and opposite bits $e^{+i0.0264089}$.
All 15 pairs commute; their phases add into the single joint SNAP table.
There are three within-cavity pairs per cavity and nine cross-cavity pairs.

The local shell counts the number $k$ of encoded 1 bits:
\begin{equation}
 f_\ell(k)=b_{\ell1}k+b_{\ell2}\binom{k}{2}+b_{\ell3}\binom{k}{3},
 \qquad S_\ell^Z=e^{-i[f_\ell(\hat k_A)+f_\ell(\hat k_B)]}.
\end{equation}
\example Fock $n=5$ is binary $101$, so $k=2$, not $5$.
For the saved first-layer coefficients $(0.8301622,-0.0460802,-0.0339120)$,
$f_1(2)=1.6142441$ radians. The shell phases for $k=0,1,2,3$ are
$0,0.8301622,1.6142441,2.3183338$. Two cavities both having $k=2$
acquire the summed phase $3.2284882$ radians.
\takeaway{The parameters $\gamma,a,b$ are shared across layers; each layer adds three shell parameters shared by both cavities: $3+3L$ parameters total.}

\clearpage
\lesson{7. How do we measure energy and spin-glass correlations?}{A histogram estimates probabilities in the basis physically measured.}
A Hadamard swaps the single-spin axes:
\begin{equation}
 \boxed{H Z H=X,\qquad H X H=Z.}
\end{equation}
\example If the original state is $\ket+$, its original $X$ value is
certainly $+1$. In the transformed state $H\ket+=\ket0$, native
computational readout returns 0, which is assigned sign $+1$.
This is why omitting the final $W$ allows native Fock readout to estimate
original-spin $X$ observables. To access original $Z$ correlations,
apply the appropriate physical basis-change pulse before number readout.

For measured bits, assign $s_i=+1$ for 0 and $s_i=-1$ for 1.
The same shot gives every single-spin sign and every pair product.
\begin{equation}
 \widehat{\langle Z_iZ_j\rangle}=\frac1M\sum_{t=1}^M s_i^{(t)}s_j^{(t)},
 \qquad
 E=\sum_{i<j}J_{ij}\langle Z_iZ_j\rangle-B\sum_i\langle X_i\rangle.
\end{equation}
\example In a correctly calibrated Z-basis measurement, suppose a pair
matches on 6 of 8 shots and differs on 2. Its correlation estimate is
$(6-2)/8=0.5$. With $J_{ij}=0.4$, this term contributes $0.2$ to energy.
An X-basis histogram similarly supplies all $X_i$ and $X_iX_j$ products,
but not the missing Z-basis probabilities.

\textbf{Why classical Hadamards cannot fix a histogram.}
$\ket{00}$ and $\ket{01}$ both have uniform probabilities for all four
outcomes when both spins are measured in X. Yet their Z-pair correlations
are $+1$ and $-1$. The X histogram alone cannot distinguish them, even
with infinitely many shots. Amplitudes and relative phases contain more
information than a histogram.

Our equal-time spin-glass indicator is
\begin{equation}
 \boxed{\chi_{\rm SG}=\frac16\sum_{i,j}\langle Z_iZ_j\rangle^2
          =1+\frac13\sum_{i<j}\langle Z_iZ_j\rangle^2.}
\end{equation}
\example For ideal $\ket+^{\otimes6}$ all distinct-spin Z correlations
vanish, so $\chi_{\rm SG}=1$. For a fixed Z-basis bit string every pair
correlation has magnitude one, so $\chi_{\rm SG}=1+15/3=6$.
Squaring a noisy correlation estimate adds sampling bias; the stored
estimator subtracts an estimated variance correction.

The 8192 shots are a total budget across settings. The retained calibrated
example uses four pairs of local settings with counts $179,1243,1291,5479$.
The precise allocation depends on $J$, $B$ and the response model.
\takeaway{Measure in the required physical bases. A finite-size correlation indicator can show a crossover; by itself N=6 does not establish a thermodynamic phase transition.}

\clearpage
\lesson{8. How does the Hamiltonian become a noisy simulation?}{The Hamiltonian gives coherent motion; Lindblad terms describe irreversible changes.}
A density matrix $\rho$ describes both pure states and statistical mixtures.
The Markovian model evolves it according to
\begin{equation}
 \boxed{\dot\rho=-i[h(t),\rho]+\sum_\mu\D[L_\mu]\rho,\qquad
 \D[L]\rho=L\rho L^\dagger-\tfrac12\{L^\dagger L,\rho\}.}
\end{equation}
The commutator is the pulse-driven evolution. Each collapse operator $L$
specifies a physical noise channel; its coefficient sets the rate.

\example With no drive, start a cavity in $\ket1$. Photon loss gives
$\rho(t)=e^{-t/T_1}\ket1\bra1+(1-e^{-t/T_1})\ket0\bra0$.
At one lifetime, $T_1=600\,\mu$s, the one-photon population is $36.8\%$
and the vacuum population is $63.2\%$. Their sum remains one.

\begin{center}
\begin{tabular}{p{0.25\linewidth}p{0.32\linewidth}p{0.32\linewidth}}
\toprule
Noise & Collapse operator & Physical meaning\\\midrule
Cavity loss & $L=\sqrt{\kappa}\,a$, $\kappa=1/T_{1c}$ & Photons leave the cavity.\\[4pt]
Transmon relaxation & $L=\sqrt{\gamma_1}\,\sigma^-$, $\gamma_1=1/T_{1q}$ & A transmon decays from $e$ to $g$.\\[4pt]
Transmon pure dephasing & $L=\sqrt{\gamma_\phi/2}\,\sigma_z$ & Phase information is randomized.\\[4pt]
Optional cavity dephasing & $L=\sqrt{\eta}\,n$ & Different Fock components lose relative coherence.\\\bottomrule
\end{tabular}
\end{center}
\example Cavity $T_1=600\,\mu$s gives $\kappa=0.0016667/\mu$s.
For an undriven $\ket7$, the first photon-loss rate is $7\kappa$, so its
mean waiting time is $600/7=85.7\,\mu$s. The jump takes it to $\ket6$,
changing encoded $111$ to $110$. This is not independent noise on three
physical spin qubits. Transmon $T_1=50\,\mu$s gives $\gamma_1=0.02/\mu$s$;$
an initially excited idle transmon survives for $10\,\mu$s with probability
$e^{-10/50}=81.9\%$. Optional cavity dephasing is set to $\eta=0$ here.

The total coherence time $T_2$ includes relaxation:
\begin{equation}
 \gamma_\phi=\frac1{T_2}-\frac1{2T_1}.
\end{equation}
\example Transmon $T_1=50\,\mu$s and $T_2=30\,\mu$s give
$\gamma_\phi=0.0233333/\mu$s and $T_\phi=42.857\,\mu$s.
The total off-diagonal coherence decays as $e^{-t/30}$, not $e^{-t/42.857}$.
\takeaway{The cavity lifetime was user supplied. The transmon lifetimes and other device parameters are assumptions, not measurements.}

\clearpage
\lesson{9. How do we simulate the whole pulse sequence?}{Use the actual drive waveforms and keep the same noise model active throughout.}
The procedure is: initialize vacuum and ground transmons; insert each pulse
bin into $h(t)$; evolve through preparation, joint SNAP, local H and local
shell; then compute the desired observable. Any physical readout pulse
must be appended if its noise is to be included.

The simulation retains 24 levels per cavity so that drives can leave the
logical levels $0,\ldots,7$. With three two-level transmons the Hilbert-space
dimension is $2^3\times24^2=4608$. A pure state has 4608 complex amplitudes;
a density matrix has $4608^2=21{,}233{,}664$ complex entries.

This motivates quantum trajectories: evolve many random pure-state
histories whose average reproduces the density matrix in the convergence
limit. Between jumps, use
\begin{equation}
 \boxed{h_{\rm eff}(t)=h(t)-\frac{i}{2}\sum_\mu L_\mu^\dagger L_\mu.}
\end{equation}
The unnormalized state's shrinking norm gives the probability of no jump.
If a jump occurs, select its channel with weight proportional to
$\|L_\mu\ket\psi\|^2$ and normalize the resulting state. The current
implementation uses finite time steps, so timestep convergence matters.

\example For one undriven photon, the no-jump amplitude is
$e^{-\kappa t/2}$ and its squared norm is $e^{-\kappa t}$.
A photon-loss jump sends $\ket1$ to $\ket0$. Averaging many such histories
reproduces the cavity mixture from Lesson 8.

\textbf{Two distinct kinds of sampling.}
\begin{itemize}
\item A trajectory samples a physical noise history inside the simulator.
For example, 1024 trajectories estimate the noisy density matrix's overlap
with a target. They are not the same as 1024 experimental measurements.
\item A measurement shot samples an outcome from the final measurement
probabilities. The 8192-shot energy budget concerns this second step.
One can simulate noise first and then sample measurement outcomes.
\end{itemize}

\textbf{What we checked.} The ideal six-spin ground state was independently
computed by diagonalizing a $64\times64$ matrix. An independent sparse,
double-precision solver verified noiseless physical pulse evolution.
For noise we compared 5 ns and 2.5 ns steps, using 1024 trajectories each
on GPU 1. Their ground-state fidelities agree within sampling uncertainty.
This is evidence of consistency at that precision, not proof of zero
finite-step error.

\textbf{What cutoff and model checks mean.} Increasing the cavity cutoff
checks numerical truncation. It does not establish that a two-level
transmon and dispersive approximation remain accurate at all populated
photon numbers. That requires calibrated couplings, detunings and
anharmonicities. Finite temperature, correlated baths and electronics
imperfections are omitted from the current baseline.
\takeaway{Quantum trajectories and a density-matrix solver describe the same chosen Lindblad model when numerical errors are controlled.}

\clearpage
\lesson{10. What do the fidelity numbers actually mean?}{Always name the target state, the output frame, and which auxiliary conditions are imposed.}
Ground-state fidelity asks how much of the final state overlaps the target.
Because the schedule stops before its final $W$, transform the target too:
\begin{equation}
 \boxed{\ket{g_X}=VW\ket{g_{\rm SK}},\qquad
 F_{\rm cav}=\bra{g_X}\rho_{\rm cav}\ket{g_X}.}
\end{equation}
Here $V$ is just the binary-to-Fock embedding from Lesson 1, and
$\rho_{\rm cav}$ is obtained by tracing out all three transmons.
\example For one spin whose target is $\ket+$, an omitted final H means the
output-frame target is $\ket0$. Comparing output $\ket0$ to $\ket+$ would
incorrectly give $1/2$; comparing it to the transformed target gives one.

We also compute a stricter overlap requiring all transmons to finish in $g$.
It is at most $F_{\rm cav}$, and neither overlap is normalized by a success
probability. For example, if the cavities are exactly on target but the
auxiliaries finish in $ggg$ only 90\% of the time, $F_{\rm cav}=1$ and
$F_{ggg}=0.9$.

\textbf{Our verified example:} $N=6$, $L=1$, $B=0.183$, instance seed
$831106000$, shorter H, $21.55\,\mu$s preparation. Final readout noise is
excluded; these are fixed ideal-optimized parameters.
\begin{center}
\begin{tabular}{lrr}
\toprule
Calculation & Cavity fidelity & Also require $ggg$\\\midrule
Ideal ansatz & $76.49\%$ & $76.49\%$\\
Noiseless pulses, independent solver & $71.844\%$ & $71.797\%$\\
Noisy, 1024 trajectories, 5 ns & $35.13\%\pm0.90$ pp & $32.70\%\pm0.92$ pp\\
Noisy, 1024 trajectories, 2.5 ns & $34.66\%\pm0.90$ pp & $32.26\%\pm0.93$ pp\\\bottomrule
\end{tabular}
\end{center}
``pp'' means percentage points. The error bars are one standard error,
$s_F/\sqrt{N_{\rm traj}}$. For example, an overlap standard deviation of
$0.288$ over 1024 trajectories gives $0.288/32=0.009$, or $0.90$ percentage
points. It does not include uncertainty in the assumed lifetimes.

The ideal ansatz is already imperfect; pulse synthesis lowers its actual
ground-state fidelity further; physical noise lowers it to about 35\% in
this instance. These effects are not independent probabilities that can
simply be multiplied. In particular, the earlier 94.7\% noiseless number
was overlap with the ideal ansatz, which is a different target.

\textbf{Where to read more.} The companion methods PDF gives full notation
and references. Useful primary papers are
\href{https://arxiv.org/abs/1503.01496}{Heeres et al. (SNAP)},
\href{https://doi.org/10.1103/PhysRevA.92.040303}{Krastanov et al. (oscillator control)},
\href{https://arxiv.org/abs/1601.05505}{Wang et al. (two cavities)},
\href{https://arxiv.org/abs/1808.07960}{Zhang et al. (driven mode exchange)}, and
\href{https://arxiv.org/abs/2005.12667}{Blais et al. (circuit-QED review)}.
For the noise equations see
\href{https://doi.org/10.1007/BF01608499}{Lindblad (1976)} and
\href{https://doi.org/10.1103/PhysRevLett.68.580}{Dalibard et al. (1992)}.
\takeaway{The verified statement is about 35\% ground-state fidelity in the chosen noisy model, for this one circuit and instance, before physical readout.}
\end{document}
